\chapter{录取之后：资格复审与入学准备}
\label{ch15}
\ist{入学前后}
\begin{keybox}
这一章处理的是“录取之后”这一段最容易被忽略的时间：从确认待录取到研究生报到，中间隔着将近一年。很多人以为拿到待录取就结束了，其实还有三道关口。第一道是入学前的资格复审，多数学校安排在录取当年6月，审的是学分是否修满、最后一学年有没有不及格、报名至报到期间有没有处分、毕业论文成绩是否达标、能不能拿到毕业证和学位证——任何一项不过，录取资格就可能被取消，这在全国有《普通高等学校学生管理规定》（教育部令第41号）和各校推免办法两套依据。第二道是报到与注册，双证原件、人事档案、政审表、体检都要在这一步到位。第三道是入学后3个月内学校按国家招生规定进行的复查，弄虚作假、徇私舞弊的取消学籍，身心状况不适宜在校学习的可保留入学资格。本章据此分五节：资格复审、休学、更换导师、提前进组、延毕风险。凡属院校特定的部分（复审材料清单、时间点、成绩与论文门槛），都点明学校与年份；凡属经验判断的部分（怎么跟导师开口、提前进组值不值、要不要换导师），归入“经验提示”，与规则条文分开。
\end{keybox}

\section{休学}
\begin{qa}{E0312}{毕业证、学位证没有按时拿到会怎么样？}
结论：后果是录取资格直接无效，几乎没有补救空间。教育部《2027年全国硕士研究生招生工作管理规定》第七十条写明，应届本科毕业生须在录取当年入学前取得国家承认的本科毕业证书，未取得的“其录取资格无效”（\nref{S0442}）。各校表述更直白：中山大学规定推免生“入学前须本科毕业并获得学士学位”，一旦查证不实或不符合要求，取消申请资格、录取资格甚至学籍（\nref{S0167}）；电子科技大学把“录取当年入学前未取得国家承认的本科毕业证”直接列为取消推免录取资格的情形（\nref{S0495}）；哈尔滨工业大学要求报到时交验身份证与学历学位证书原件，报到前未取得毕业证即取消录取资格（\nref{S0066}）。

注意学位证与毕业证是两件事：多所学校同时要求“获得学士学位”（\nref{S0635}、\nref{S0604}），只拿到毕业证不等于过关。
\nref{S0442}
\nref{S0167}
\nref{S0495}
\nref{S0066}
\nref{S0635}
\nref{S0604}
\end{qa}

\iin{教育部学生服务与素质发展中心（学信网/研招网）}
\ith{证书要求}
\ist{入学前后}
\begin{qa}{E0316}{政审表怎么填？由谁来签字盖章？}
结论：表由你取回、由你所在单位填，但你本人不得经手，必须密封转递。浙江大学能源工程学院2026年通知规定，政审表交本人档案所在高校二级学院如实填报，需填写完整、含所在单位负责人签字及单位公章方为有效，“考生所在单位政审意见”须加盖二级院系党组织公章；并明确政审表应由政审单位寄送、密封且在封口处加盖单位公章，“考生本人严禁接触政审表”（\nref{S0542}）。教育部文件也要求，人事档案审查或政审必须在发放录取通知书之前完成（\nref{S0046}）。中山大学的做法是把“现实表现复审表”放进研招系统由拟录取考生自行下载，再按调档函要求寄回（\nref{S0164}）；哈尔滨工业大学要求现实表现材料由本人档案或工作所在单位的人事、政工部门加盖印章（\nref{S0066}）。

\begin{itemize}
  \item 谁填：本科院系辅导员与二级党组织
  \item 谁签章：单位负责人签字 + 二级院（系）党组织公章
  \item 谁寄：由政审单位寄出并密封、封口盖章
\end{itemize}
\nref{S0542}
\nref{S0046}
\nref{S0164}
\nref{S0066}
\end{qa}

\iin{浙江大学}
\ith{政审}
\ist{入学前后}
\begin{qa}{E0318}{档案怎么转到研究生学校？自己带过去行不行？}
结论：走“调档函 + 公对公转递”，个人自带通常不被接受。标准流程由招生单位发起：拟录取名单确定后，招生单位向考生所在单位函调人事档案（或档案审查意见）和本人现实表现等材料，函调材料须由本人档案或工作所在单位的人事、政工部门加盖印章（\nref{S0051}、\nref{S0066}）。中山大学把调档函放在研招系统里，由拟录取硕士（含免试生）自行下载，并按调档函要求将材料按时寄回（\nref{S0164}）；浙江大学能源工程学院明确调档一般安排在入学当年5月左右，调档函通过电子邮件发放（\nref{S0542}）；哈尔滨工业大学要求调档工作在10月底前完成（\nref{S0066}）。

注意：推免生同样要调档，不会因为免初试而跳过这一步。
\nref{S0051}
\nref{S0066}
\nref{S0164}
\nref{S0542}
\end{qa}

\iin{中华人民共和国教育部}
\ith{档案转递}
\ist{入学前后}
\begin{qa}{E0319}{档案里缺材料（比如入党材料、成绩单）怎么办？}
结论：档案缺件不是“取消资格”的硬条款问题，而是会卡住调档、政审和入学复查的流程问题。国家层面没有“档案缺某材料即取消入学资格”的统一规定；现有的依据是两方面：一是调档与政审要求函调的人事档案（或档案审查意见）须由档案所在单位人事、政工部门加盖印章（\nref{S0051}、\nref{S0066}），材料不齐时对方往往无法出具审查意见；二是入学后学校要做全面复查，西北工业大学明确把“档案复查”列入入学后3个月内的复查内容（\nref{S0506}）。也就是说，缺件的后果是流程被卡住，而不是被直接取消。

\begin{itemize}
  \item 入党材料：回本科院系党组织补齐
  \item 成绩单：教务处重新打印并加盖公章
  \item 登记表类：联系原填表单位补办
\end{itemize}
\nref{S0051}
\nref{S0066}
\nref{S0506}
\end{qa}

\iin{中华人民共和国教育部}
\ith{档案缺件}
\ist{入学前后}
\begin{qa}{E0322}{党组织关系怎么转？介绍信抬头写什么？}
结论：原则是入学前把组织关系转入录取学校，但抬头怎么写属各校规定，必须问录取院系。现有材料中可确认的是学校要求转入：安徽工业大学博士招生简章要求全日制非定向考生“入学前须将全部人事档案、组织关系等转入我校”（\nref{S0161}）；中国科学技术大学与宁波东方理工大学的联合培养博士生，由中科大“按相关规定转接人事档案、组织关系等”（\nref{S0362}）。《普通高等学校学生管理规定》只对学籍、档案作出安排（\nref{S0040}），其中并不涉及党员组织关系转接。

\begin{itemize}
  \item 向谁问：录取院系党组织或研究生辅导员
  \item 要什么：准确的一级党组织抬头 + 转接方式（线上系统或纸质介绍信）
  \item 何时办：一般在收到录取通知后、开学前
\end{itemize}
\nref{S0161}
\nref{S0362}
\end{qa}

\iin{安徽工业大学}
\ith{党组织关系}
\ist{入学前后}
\begin{qa}{E0323}{预备党员入学后怎么转正？会不会影响？}
结论：跨校本身不影响转正，真正的风险是培养考察材料断档。现有材料中未见针对推免生预备党员转正的专门条款。可依据的关联规则是：《普通高等学校学生管理规定》第十一条规定，学生入学后学校应当在3个月内进行复查，复查内容包括本人及身份证明与录取通知、考生档案等是否一致（\nref{S0040}）；而调档环节要求人事档案与审查意见由档案所在单位盖章出具（\nref{S0051}、\nref{S0066}）。党员材料是人事档案的组成部分，材料不全就可能在转接与复查环节卡住。

\begin{itemize}
  \item 需衔接：培养考察登记表、党课结业证明、思想汇报、预备期起止时间
  \item 常见做法：由本科党支部密封装入档案，入学后第一时间向新支部报到
  \item 时间点：预备期满前提出转正申请
\end{itemize}
\nref{S0040}
\nref{S0051}
\nref{S0066}
\end{qa}

\iin{中华人民共和国教育部}
\ith{党员转正}
\ist{入学前后}
\begin{qa}{E0324}{智慧团建上的团组织关系怎么转？}
结论：团关系转接是入学手续的一环，但“智慧团建”的具体操作没有全国统一条文可引。现有材料只能确认学校要求转入：安徽工业大学要求全日制非定向考生入学前将人事档案、组织关系等转入我校（\nref{S0161}）；中国科学技术大学与宁波东方理工大学的联合培养博士生，由中科大统一转接人事档案、组织关系（\nref{S0362}）。也就是说，转是必须转的，但选哪个团支部、何时发起，由录取院系团委规定。

\begin{itemize}
  \item 通常路径：录取院系团委在智慧团建中建立新生团支部，你发起转入申请
  \item 需要：团员证、入团志愿书（一般在档案内）
  \item 时间：多在开学后由班级团支书统一组织
\end{itemize}
\nref{S0161}
\nref{S0362}
\end{qa}

\iin{安徽工业大学}
\ith{团组织关系}
\ist{入学前后}
\begin{qa}{E0326}{入学体检和复试体检有什么区别？}
结论：区别在“谁在什么时候组织、不合格的后果是什么”。教育部规定，考生体检工作由招生单位在考生拟录取后组织进行，各招生单位参照《普通高等学校招生体检工作指导意见》并结合专业实际提出本单位体检要求（\nref{S0442} 第七十二条）。所以“复试体检”通常就是拟录取后的这一次，与录取结果直接挂钩：浙江大学能源工程学院要求拟录取考生自行到校医院或当地二甲以上医院体检，“体检不合格者取消拟录取资格”（\nref{S0542}）。入学后的体检则与学籍绑定：哈尔滨工业大学规定已录取推免生体检在入学报到时进行（\nref{S0066}）；河北大学规定体检在入学报到后进行，体检不合格者经研究生院批准休学，休学期满后仍不合格者取消学籍（\nref{S0604}）。

\par\vspace{2pt}\nobreak\begin{center}\footnotesize
\begin{tabularx}{\textwidth}{>{\raggedright\arraybackslash}X>{\raggedright\arraybackslash}X>{\raggedright\arraybackslash}X}
\toprule
\textbf{项目} & \textbf{拟录取/复试体检} & \textbf{入学报到体检} \\ \midrule
组织者 & 招生单位（拟录取后） & 录取学院/校医院 \\
后果 & 不合格取消拟录取资格 & 不合格先休学，期满仍不合格取消学籍 \\
\bottomrule
\end{tabularx}
\end{center}\vspace{2pt}\par
\nref{S0442}
\nref{S0542}
\nref{S0066}
\nref{S0604}
\end{qa}

\iin{教育部学生服务与素质发展中心（学信网/研招网）}
\ith{体检差异}
\ist{入学前后}
\begin{qa}{E0338}{保留入学资格、休学怎么申请？需要什么理由？}
结论：这是两条不同路径——保留入学资格发生在入学前，休学发生在取得学籍之后，条件与期限都由学校规定。教育部令第41号（2017年施行）给出全国统一的框架（\nref{S0040}）：

\begin{itemize}
  \item 保留入学资格：新生可以申请，保留期间不具有学籍；期满前应向学校申请入学，经审查合格后办理入学手续，审查不合格的取消入学资格，逾期未办且无不可抗力等正当理由的视为放弃入学资格（第十条）。教育部招生规定同样授权被录取新生可申请保留入学资格，条件期限由招生单位规定，并纳入录取当年招生计划管理（\nref{S0442} 第七十一条）。
  \item 休学：学生申请或学校认为应当休学的，经学校批准可以休学，休学次数和期限由学校规定；休学期间学校保留学籍但不享受在校学生待遇（第二十五、二十八条）。
  \item 应征入伍：新生和在校生应征参加中国人民解放军（含武警），学校应当保留其入学资格或学籍至退役后2年（第二十七条）。
  \item 复学：休学期满前应在学校规定期限内提出复学申请，经学校复查合格方可复学（第二十九条）。
\end{itemize}
\nref{S0040}
\nref{S0442}
\end{qa}

\iin{中华人民共和国教育部}
\ith{休学}
\ist{入学前后}
\begin{qa}{E0345}{宿舍什么时候申请？不分宿舍怎么办？}
结论：多数学校为全日制研究生提供住宿，但要你主动申请、费用自理，具体时间在录取后才通知。厦门大学2026年招生简章写明“我校为正常学制年限内的全日制研究生提供住宿，由学生提出申请，住宿费用自理”，非全日制不予安排；硕士新生为2—6人/间，住宿费800—1600元/人·学年，“具体住宿安排在录取后另行通知”（\nref{S0529}）。电子科技大学为全日制非定向硕、博士研究生提供住宿，博士生两人一间（\nref{S0495}）。

\begin{itemize}
  \item 关注渠道：研究生院“研究生管理/后勤”栏目、录取通知书随附的入学须知、院系通知群
  \item 常见时点：入学须知发布后至开学前后，部分学校按院系分片、先到先得
  \item 不安排住宿时：需自行在校外租房，入学前先了解校区周边行情
\end{itemize}
\nref{S0529}
\nref{S0495}
\end{qa}

\iin{厦门大学}
\ith{宿舍申请}
\ist{入学前后}
\begin{qa}{E0346}{入学英语分级考试、免修是怎么回事？}
结论：分级与免修属院校自主安排，国家层面没有统一规定。《普通高等学校学生管理规定》（\nref{S0040}）只作授权性规定：考核分为考试和考查两种，考核和成绩评定方式以及考核不合格的课程是否重修或补考，由学校规定（第十三条）；学生根据学校有关规定可以申请修读课程，修读的课程成绩（学分）经学校审核同意后予以承认（第十六条）。换句话说，分不分级、能不能免修、免修依据是四六级还是入学测试，全看你所在学校的研究生英语教学办法，教育部文件中查不到统一条款。

\begin{itemize}
  \item 常见形式：入学后校内英语水平测试分级，按级别修读相应课程
  \item 常见免修依据：六级达到一定分数、考研英语高分、雅思托福成绩
  \item 免修与免考是两回事，多数学校“免修不免考”
\end{itemize}
\nref{S0040}
\end{qa}

\iin{中华人民共和国教育部}
\ith{英语分级}
\ist{入学前后}
\begin{qa}{E0347}{推免生入学后能不能申请硕博连读？}
结论：可以，但走的是校内二次选拔，属院校规定，需要提前规划成绩与导师名额。硕博连读是博士招生的独立方式之一：中国农业大学把博士招生分为直博生、硕博连读、申请考核三种，其中硕博连读是“选拔校内优秀在读硕士生”，由各招生院系依据学校的硕博连读选拔办法组织实施，博士阶段学制3年（\nref{S0160}）。安徽工业大学的硕博连读面向本校全日制非定向在读硕士（含学术学位和专业学位），本科须获学士学位，博士阶段学制3年（\nref{S0161}）。天津医科大学要求面向二、三年级学术学位非定向硕士，须修完全部学位课程、无不及格或补考且课程成绩优良（平均80分以上，公共英语不低于80分）（\nref{S0184}）。

三点差异：申请时点、成绩门槛、是否限同一导师，各校不同。
\nref{S0160}
\nref{S0161}
\nref{S0184}\par\textbf{本版实操补充：} 是否能申请取决于学校的硕博连读制度、学科名额和考核，不会因推免录取身份自动取得博士资格。入学前先查培养方案对申请年级、课程成绩、导师同意、资格考试和退出方式的规定；如果当前主要目标是读硕士，选择导师和项目时也应确认硕士阶段培养安排。硕博连读属于后续培养路径，不能把招生宣传中的“有机会”理解成保证转博。
\end{qa}

\iin{中国农业大学}
\ith{硕博连读}
\ist{入学前后}
\section{更换导师}
\begin{qa}{E0313}{大四最后一学期挂科了，会不会被取消录取？}
结论：会，这是资格复审里最常被卡的一条。清华大学的资格复审要求明确“本科最后一学年课程考核成绩不得出现不及格记录”，未通过资格复审者取消研究生录取资格（\nref{S0491}、\nref{S0492}、\nref{S0254}）。西华师范大学外国语学院把“本科学习阶段第7-8学期学习成绩有不及格科目”直接列为取消录取资格的情形（\nref{S0152}）。河北大学2027年办法规定，对未达到推荐高校提出的本科应届毕业学年学业成绩等要求的推免生，由推荐高校取消其推免生资格并通报，我校取消其推免录取资格（\nref{S0604}）；中山大学也有拟录取后至本科毕业前发生相关规定情形即取消推免录取资格的表述（\nref{S0167}）。

注意认定口径：多数学校看的是成绩单上是否“出现不及格记录”，补考及格后原始记录是否留存，各校不同。
\nref{S0491}
\nref{S0492}
\nref{S0254}
\nref{S0152}
\nref{S0604}
\nref{S0167}
\end{qa}

\iin{清华大学}
\ith{学业变化}
\ist{入学前后}
\begin{qa}{E0314}{毕业论文没过会不会影响研究生入学？}
结论：会，因为毕业论文与毕业证学位证直接绑定，而双证是录取生效的前提。清华大学的资格复审要求“本科综合论文训练（毕业设计）应为优良或80分（含）以上”，未通过者取消录取资格（\nref{S0491}、\nref{S0254}）；西华师范大学外国语学院要求“毕业设计（论文）未取得良好以上成绩”者取消录取资格（\nref{S0152}）。国家层面是链条式的：学生在规定学习年限内修完教学内容、成绩合格、达到毕业要求才准予毕业并发放毕业证书，未达到毕业要求的只能结业（\nref{S0040} 第三十二、三十三条）；而应届本科生须在录取当年入学前取得本科毕业证书，未取得的其录取资格无效（\nref{S0442} 第七十条）。

论文不及格 → 拿不到毕业证 → 录取资格无效，这条链没有断点。
\nref{S0491}
\nref{S0254}
\nref{S0152}
\nref{S0040}
\nref{S0442}
\end{qa}

\iin{清华大学}
\ith{毕业论文}
\ist{入学前后}
\begin{qa}{E0321}{农村户口迁到学校后还能迁回农村吗？}
结论：未见全国统一规定，能否迁回取决于生源地户籍政策，且通常不可逆。《普通高等学校学生管理规定》涉及的唯一户口条款是退学情形：“退学学生的档案由学校退回其家庭所在地，户口应当按照国家相关规定迁回原户籍地或者家庭户籍所在地”（\nref{S0040} 第三十一条）。请注意条文说的是“原户籍地”，而不是“恢复农业户口性质”——这是两个概念。除此之外，教育部令第41号没有对在校生户口迁移作出任何规定，也没有承诺迁出后可恢复原户口性质。所以“能不能迁回农村”不是教育系统能回答的问题，要看生源地公安机关的现行户籍政策。
\nref{S0040}
\end{qa}

\iin{中华人民共和国教育部}
\ith{户口回迁}
\ist{入学前后}
\begin{qa}{E0325}{入学体检复检不合格会不会取消录取？}
结论：可能，但处理顺序通常是“先休学、再取消学籍”，不是一票出局。教育部规定体检由招生单位在拟录取后组织，参照《普通高等学校招生体检工作指导意见》并结合专业实际提出本单位要求（\nref{S0442} 第七十二条）；浙江大学能源工程学院明确体检不合格者取消拟录取资格（\nref{S0542}）。入学后的处理更有层次：河北大学2027年办法规定体检不合格者经研究生院批准休学，休学期满后体检仍不合格者取消学籍（\nref{S0604}）；《普通高等学校学生管理规定》第十一条也规定，复查中发现学生身心状况不适宜在校学习、经学校指定二级甲等以上医院诊断需要在家休养的，可以保留入学资格（\nref{S0040}）。西北工业大学则把入学后3个月内全面复查（含复查不合格按相关要求处理、情节严重移交有关部门）作为常规流程（\nref{S0506}）。
\nref{S0442}
\nref{S0542}
\nref{S0604}
\nref{S0040}
\nref{S0506}
\end{qa}

\iin{教育部学生服务与素质发展中心（学信网/研招网）}
\ith{体检复检}
\ist{入学前后}
\begin{qa}{E0327}{乙肝携带、结核病史、心理疾病会不会被退学？}
结论：三者都不构成“一律退学”，判断标准是能否保证在校正常学习生活。教育部已明确普通高等学校招生入学身体检查取消乙肝项目检测，参见《2027年全国硕士研究生招生工作管理规定》第七十二条对教学厅〔2010〕2号文件的引用（\nref{S0442}）。《普通高等学校招生体检工作指导意见》规定：慢性肝炎病人并且肝功能不正常者（肝炎病原携带者但肝功能正常者除外）、严重精神病未治愈者、结核病除已硬结稳定或治愈后满一年以上未复发等特定情形外，学校可以不予录取（\nref{S0039}）。注意条文用的是“可以”而非“必须”，且该指导意见明确只是指导性意见，高校可提出补充规定（\nref{S0039} 通知说明部分）。

\begin{itemize}
  \item 乙肝：招生入学体检已取消该项目检测
  \item 结核：已治愈且符合《指导意见》列举情形的，不属于“可以不予录取”
  \item 心理疾病：41号令第十一条给的是保留入学资格通道，不是退学通道（\nref{S0040}）
\end{itemize}
\nref{S0442}
\nref{S0039}
\nref{S0040}
\end{qa}

\iin{教育部学生服务与素质发展中心（学信网/研招网）}
\ith{健康限制}
\ist{入学前后}
\begin{qa}{E0330}{提前进组发现导师不合适，还能不能换？}
结论：能换，但窗口期很短，且没有全国统一规则可依。现有材料中找不到“更换导师”的全国性条款，能确认的是学校会明确师生双选的时间点：北京航空航天大学仪器科学与光电工程学院规定，推免攻读2027年研究生（博士）确认拟录取后，2027年6月前由学院根据导师名额分配情况组织双向选择导师；硕士研究生则在复试结束后统一组织师生双选（\nref{S0421}）。南京理工大学规定“推免生在拟录取后即可与导师进行互选，确定导师”（\nref{S0560}）。这些时点就是你的机会窗口。
\nref{S0421}
\nref{S0560}
\end{qa}

\iin{北京航空航天大学}
\ith{更换导师}
\ist{入学前后}
\begin{qa}{E0331}{开学前想换导师，具体怎么操作？}
结论：先查学院的双选机制什么时候关闭，再在关闭前提出。规范做法是走学院组织的师生双选：北京航空航天大学仪器科学与光电工程学院在复试结束后统一组织硕士师生双选，直博生在2027年6月前由学院按导师名额分配组织双向选择（\nref{S0421}）；南京理工大学推免生拟录取后即可与导师互选（\nref{S0560}）。只要双选流程还在进行，你重新选择一位有招生名额的导师即可；一旦学院公示或系统锁定名单，就要走学院层面的变更流程了。

\begin{itemize}
  \item 找谁：研究生教学秘书或辅导员，问清书面流程与截止日
  \item 前提：目标导师有招生名额并同意接收
  \item 顺序：先确认接收，再报备学院，最后才与原导师沟通
\end{itemize}
\nref{S0421}
\nref{S0560}
\end{qa}

\iin{北京航空航天大学}
\ith{更换导师}
\ist{入学前后}
\begin{qa}{E0332}{换导师需要原导师同意吗？会不会撕破脸？}
结论：没有统一条款规定必须原导师书面同意，但现实里绕不过他——因为换导师的本质是换名额。现有材料中未见“更换导师须原导师同意”的全国性或高校明文规定。能确认的是导师名额的硬约束：北京航空航天大学规定直博生结合考生意向导师以及导师的招生指标从高到低拟录取，“若考生因意向导师招生指标等原因未被其拟录取，且考生不同意调整导师，则视为放弃拟录取”（\nref{S0421}）。这句话反过来说明，导师与名额是绑定的，换导师必须有一个愿意接收你且还有指标的新导师。

\begin{itemize}
  \item 学院流程上需要的通常是新导师接收意见 + 学院审批
  \item 原导师的“同意”更多是事实上的，而非文件上的
\end{itemize}
\nref{S0421}
\end{qa}

\iin{北京航空航天大学}
\ith{更换导师}
\ist{入学前后}
\begin{qa}{E0339}{入学后能不能转专业或者换方向？}
结论：可以申请，但研究生阶段受限更多，且以特殊招生形式录取的不得转。《普通高等学校学生管理规定》第二十一条规定（\nref{S0040}），学生在学习期间对其他专业有兴趣和专长的，可以申请转专业；但“以特殊招生形式录取的学生，国家有相关规定或者录取前与学校有明确约定的，不得转专业”，具体办法由学校制定并须公示。转学另设门槛：研究生拟转入学校、专业的录取控制标准高于其所在学校、专业的不予转学，且“研究生转学还应当经拟转入专业导师同意”（第二十二、二十三条）。电子科技大学也明确，研究生在学习过程中不再进行专业学位类别或专业领域的选择（\nref{S0495}）。

推免生是按专业录取的，是否属于“录取前有明确约定”，各校口径不同，需先问清楚。
\nref{S0040}
\nref{S0495}
\end{qa}

\iin{中华人民共和国教育部}
\ith{转专业}
\ist{入学前后}
\begin{qa}{E0349}{能不能提前毕业或者缩短学制？}
结论：可以申请提前毕业，但条件由学校定，且有“最长学习年限”的天花板。《普通高等学校学生管理规定》第三十二条规定（\nref{S0040}），学生提前完成教育教学计划规定内容、获得毕业所要求的学分，可以申请提前毕业，提前毕业的条件由学校规定；同时第二十五条规定，除另有规定外，学生应当在学校规定的最长学习年限（含休学和保留学籍）内完成学业。学费也按实际学习年限计算：电子科技大学依据川发改价格〔2024〕417号，对提前毕业或延期毕业的学生按实际学习年限收费（\nref{S0495}）。

\begin{itemize}
  \item 决定因素：学位论文进度 \textgreater{} 学分完成情况
  \item 各校学制不同（硕士多为2—3年），提前毕业门槛也不同
  \item 提前毕业通常不免修培养方案规定的环节
\end{itemize}
\nref{S0040}
\nref{S0495}
\end{qa}

\iin{中华人民共和国教育部}
\ith{提前毕业}
\ist{入学前后}
\section{提前进组}
\begin{qa}{E0317}{要不要开无犯罪记录证明？去哪开？}
结论：研究生入学没有全国统一的“无犯罪记录证明”要求，真正规范的是思想政治素质和品德考核与政审。教育部规定，拟录取名单确定后，招生单位应向考生所在单位函调人事档案（或档案审查意见）和本人现实表现等材料，全面考查思想政治和品德情况，考核不通过者不予录取（\nref{S0051}）；人事档案审查或政审必须在发放录取通知书之前完成（\nref{S0046}）。推免遴选层面则常把“无违法犯罪记录”直接写成推荐条件：成都大学商学院要求品行表现优良、无违法犯罪记录和尚未解除的纪律处分，思想品德考核不合格者不予推荐（\nref{S0079}）。

所以：除非录取学校明确要求，一般不需要单独开具。
\nref{S0051}
\nref{S0046}
\nref{S0079}
\end{qa}

\iin{中华人民共和国教育部}
\ith{证明材料}
\ist{入学前后}
\begin{qa}{E0320}{户口要不要迁到学校？迁了有什么好处？}
结论：属自愿选择，没有统一规定，迁与不迁各有利弊。《普通高等学校学生管理规定》唯一涉及户口的是退学情形：退学学生的档案由学校退回其家庭所在地，“户口应当按照国家相关规定迁回原户籍地或者家庭户籍所在地”（\nref{S0040} 第三十一条）。除此之外，现行教育部规章并未对在校生户口迁移作出规定，所以“要不要迁”不是教育系统的问题，而是学校、地方户籍政策和个人情况的综合判断。

\begin{itemize}
  \item 迁的好处：在校期间办证、出国、驾照、购房摇号等按当地口径可能更方便（各城市政策不同）
  \item 不迁的好处：保留原户籍地绑定的权益，家里的事不受影响
  \item 关键变量：学校是否强制、迁入后能否落回、落户地址是学校集体户还是街道
\end{itemize}
\nref{S0040}
\end{qa}

\iin{中华人民共和国教育部}
\ith{户口迁移}
\ist{入学前后}
\begin{qa}{E0328}{录取后想提前进组，怎么跟导师开口比较合适？}
结论：核心是“请示具体安排”，而不是“我想来”。现有材料里没有提前进组的统一规则，但学校一般有正式的师生双选与培养安排节点：北京航空航天大学仪器学院在复试结束后统一组织硕士师生双选，直博生在2027年6月前完成双选（\nref{S0421}）；南京理工大学推免生拟录取后即可与导师互选（\nref{S0560}）。也就是说，通常要等双选完成或导师明确表示接收之后，“进组”这个动作才谈得上。

一封容易被答应的消息包含三件事：
\begin{itemize}
  \item 你能投入什么时间、能待多久
  \item 你能做什么（读文献、复现实验、跟师兄做小任务）
  \item 你需要什么（工位、账号、住宿建议）
\end{itemize}
\nref{S0421}
\nref{S0560}
\end{qa}

\iin{北京航空航天大学}
\ith{提前进组}
\ist{入学前后}
\begin{qa}{E0329}{提前进组有补贴吗？暑假要不要留下来？}
结论：提前进组期间的补贴没有统一规定，取决于课题组经费与学院安排；能依据的是入学取得学籍之后的奖助政策。南京理工大学明确推免录取的硕士第一学年享受学校支持性学业奖学金（\nref{S0560}）；电子科技大学列明研究生奖助学金包括国家奖学金、学业奖学金、国家助学金、“三助”助学金及各类荣誉奖励（\nref{S0495}）。这些都是入学后的待遇。

\begin{itemize}
  \item 提前进组的补贴：各课题组自定，可能给劳务、可能只解决工位与住宿
  \item 暑假留不留：先看本科毕业论文答辩与毕业手续的时间安排
  \item 若留下：确认住宿怎么解决（宿舍通常尚未开放）
\end{itemize}
\nref{S0560}
\nref{S0495}
\end{qa}

\iin{南京理工大学}
\ith{提前进组}
\ist{入学前后}
\begin{qa}{E0333}{导师没经费、快退休了，要不要换？}
结论：这是没有条文可依的判断题，只能按可核查的事实来定。现有材料能提供的间接依据有两类：一是导师名额是拟录取的硬约束——北京航空航天大学规定直博生结合考生意向导师以及导师招生指标拟录取，考生不同意的视为放弃拟录取（\nref{S0421}）；二是学制与最长学习年限由学校规定，中国农业大学最长在校学习年限为学制加2年（\nref{S0160}），河北大学直博生学制5年、最长学习年限不超过8年（\nref{S0604}）。年限越紧，选错导师的代价越大。

建议的核查顺序：当年是否还有招生资格与名额 → 在读师兄的真实状态 → 经费与年龄。
\nref{S0421}
\nref{S0160}
\nref{S0604}
\end{qa}

\iin{北京航空航天大学}
\ith{导师质量}
\ist{入学前后}
\begin{qa}{E0334}{联培、双导师制怎么确定主导师？}
结论：看学籍注册在哪所学校，主导师通常跟着学籍走。中国科学技术大学与宁波东方理工大学的联合培养博士生办法写得很明确：联合培养博士生“在中国科学技术大学注册学籍，由双方联合培养，双导师指导”，在中科大进行课程学习、在宁波东方理工从事与学位论文相关的科研工作，并由中科大按相关规定转接人事档案、组织关系等；符合中科大毕业条件的颁发中科大毕业证书，符合学位授予要求的授予中科大学位（\nref{S0362}）。教育部层面则有《高等学校和科研机构开展联合培养博士研究生工作暂行办法》作为制度框架（\nref{S0034}）。

\begin{itemize}
  \item 学籍归属 = 主导师所在单位 = 证书颁发单位
  \item 双导师制下另一位导师负责科研现场指导
\end{itemize}
\nref{S0362}
\nref{S0034}
\end{qa}

\iin{与宁波东方理工大学}
\ith{双导师}
\ist{入学前后}
\begin{qa}{E0340}{录取后学校改名、专业调整了怎么办？}
结论：机构与学科调整会传导到培养安排，但有制度兜底。国家层面正推进学科专业调整：《高等教育学科专业设置调整优化行动方案（2025—2027年）》由中央教育工作领导小组印发（\nref{S0461}）。学生侧的保护条款在《普通高等学校学生管理规定》第二十一条（\nref{S0040}）：学校根据社会对人才需求情况的发展变化需要适当调整专业的，应当允许在读学生转到其他相关专业就读。更严重的情形是第二十二条：学生因学校培养条件改变等非本人原因需要转学的，学校应当出具证明，由所在地省级教育行政部门协调转学到同层次学校。

\begin{itemize}
  \item 改名：通常只影响证书与公章，对培养无实质影响
  \item 学院合并/专业停招：可能改变培养方案与导师团队
  \item 学位点调整：影响授予学位的学科口径
\end{itemize}
\nref{S0461}
\nref{S0040}
\end{qa}

\iin{教育部学生服务与素质发展中心（学信网/研招网）}
\ith{政策变化}
\ist{入学前后}
\begin{qa}{E0341}{录取后导师跳槽了怎么办？}
结论：没有专门的全国条款，但学校有义务安排你的培养，现实路径是在同院系内转由其他导师接手。可依据的制度背景是导师在培养中的角色：教育部办公厅要求培养单位强化导师责任，“导师是培养质量第一责任人”（\nref{S0432}）；《普通高等学校学生管理规定》也把导师意见写入转学条款——研究生转学应当经拟转入专业导师同意（\nref{S0040} 第二十三条）。这说明导师变动属于院系必须处理的培养事项，而不是你个人的麻烦。

\begin{itemize}
  \item 找谁：研究生教学秘书或分管研究生工作的副院长
  \item 问什么：同方向还有哪些导师有名额、你的课题能否延续
  \item 留什么：书面确认新导师与培养方案不变
\end{itemize}
\nref{S0432}
\nref{S0040}
\end{qa}

\iin{教育部学生服务与素质发展中心（学信网/研招网）}
\ith{导师变动}
\ist{入学前后}
\begin{qa}{E0344}{录取通知书丢了怎么办？还能报到吗？}
结论：未见统一规定，但补办或凭系统记录报到在实践中可行，关键动作是第一时间联系研究生招生办公室。可参照的规则是：按国家招生规定录取的新生，持录取通知书按学校有关要求和规定的期限到校办理入学手续；因故不能按期入学的应当向学校请假，未请假或请假逾期的，除不可抗力等正当事由外视为放弃入学资格（\nref{S0040} 第八条）。而学籍注册的依据是学校对入学资格的审查结果，而不是那张纸本身——报到时学校对新生入学资格进行初步审查，审查合格即办理入学手续、予以注册学籍（第九条）。
\nref{S0040}
\nref{S0066}

先联系录取单位研招办或新生接待部门，说明姓名、身份证明、专业和邮寄信息，请其确认补寄、补开证明或报到替代材料；不要自行假设电子截图一定可代替原件。同步核实报到日期、请假渠道和是否需要在新生系统登记。学校未明确前保留招生系统录取状态、快递查询结果和沟通回执；若无法按期到校，按校规及时办理请假，避免失联导致被视为放弃入学。
\end{qa}

\iin{中华人民共和国教育部}
\ith{通知书丢失}
\ist{入学前后}
\section{资格复审}
\begin{qa}{E0310}{推免录取后资格复审要交什么材料？}
结论：材料清单各校不同，但审查维度高度一致。清华大学的资格复审要求可作为模板：已完成本科培养方案规定的所有课程及实践环节的学分要求；本科最后一学年课程考核成绩不得出现不及格记录；自报名至入学报到之日未受任何处分；本科综合论文训练（毕业设计）应为优良或80分（含）以上；入学前获得本科毕业证书和学士学位证书；符合各院系自定的不低于学校复审要求的标准（\nref{S0491}、\nref{S0492}、\nref{S0254}）。河海大学则把思想政治素质和品德考核与资格复审合并，明确四类取消录取资格情形（\nref{S0178}）。

\begin{itemize}
  \item 常见材料：本科成绩单与排名证明、现实表现复审/政审表、体检表、调档函回执
  \item 毕业证学位证：多数学校在报到时交验原件（\nref{S0066}）
  \item 常见时点：录取当年上半年，多数集中在6月前后
\end{itemize}
\nref{S0491}
\nref{S0492}
\nref{S0254}
\nref{S0178}
\nref{S0066}
\end{qa}

\iin{清华大学}
\ith{资格复审}
\ist{入学前后}
\begin{qa}{E0311}{资格复审是什么时候？开学前还是开学后？}
结论：主体在开学前完成，入学后还有一次3个月的复查，这是两件事。清华大学明确“我校将在录取当年6月对所有拟录取的推免生进行资格复审”，通过复审并经调档案审查合格后才正式录取、发放录取通知书（\nref{S0491}、\nref{S0492}）。河海大学把资格复审放在“正式发放录取通知书之前”，录取通知书在2027年6月左右统一发放（\nref{S0178}）。哈尔滨工业大学把调档安排在10月底前完成，录取通知书7—8月发放，体检放在入学报到时（\nref{S0066}）。开学后那一关是教育部统一规定的：学生入学后，学校应当在3个月内按照国家招生规定进行复查（\nref{S0040} 第十一条），复查内容包括录取资格真实性、身份与录取通知及考生档案是否一致、身心健康状况是否符合报考专业要求等；西北工业大学也列明入学后3个月内做照片、证件、档案、资格、户籍等复查（\nref{S0506}）。
\nref{S0491}
\nref{S0492}
\nref{S0178}
\nref{S0066}
\nref{S0040}
\nref{S0506}
\end{qa}

\iin{清华大学}
\ith{资格复审}
\ist{入学前后}
\begin{qa}{E0315}{资格复审时发现成绩排名掉出前多少名了，怎么办？}
结论：多数学校不设“掉出前多少名”的固定线，而是设“必须达到推荐高校的学业成绩要求”这条兜底线，一旦不达标由推荐高校取消推免资格。河北大学2027年办法写明：对于未达到推荐高校提出的本科应届毕业学年学业成绩等要求的推免生，推荐高校应取消其推免生资格并通报我校，由我校取消其推免录取资格（\nref{S0604}）。浙江大学医学院2027年综合考核细则有几乎相同的表述，并规定在我校规定的研究生新生入学日期前无法提供毕业证书和学士学位证书的将被取消入学资格（\nref{S0029}）。清华的口径则是“本科最后一学年课程考核成绩不得出现不及格记录”（\nref{S0491}）；中山大学也规定，一经查证不符合招生要求即取消申请资格、录取资格、入学资格（\nref{S0167}）。也就是说，决定权在推荐高校的推免实施办法，而不是接收学校的排名线。
\nref{S0604}
\nref{S0029}
\nref{S0491}
\nref{S0167}
\end{qa}

\iin{河北大学}
\ith{排名变化}
\ist{入学前后}
\begin{qa}{E0342}{录取后不去报到会怎样？会不会影响以后考研？}
结论：会被认定为放弃入学资格，并可能被通报不诚信，但对以后考研没有全国性的“禁考”条款。按国家招生规定录取的新生，应持录取通知书按学校要求和期限到校办理入学手续，未请假或请假逾期的，除不可抗力等正当事由外“视为放弃入学资格”（\nref{S0040} 第八条）。学校层面还有更具体的后果：复旦大学规定本专项考生确认接受拟录取结果后不得放弃，否则可取消其在校其他项目的拟录取资格并向其本科高校通报不诚信行为（\nref{S0189}）；东北师范大学规定“如因个人原因放弃录取资格，我校不予出具任何证明和手续”（\nref{S0607}）；郑州大学要求被录取的考生保证录取后能正常进校就读，不得另行就业、出国或再次报考全国研究生入学考试（\nref{S0635}）。
\nref{S0040}
\nref{S0189}
\nref{S0607}
\nref{S0635}
\end{qa}

\iin{中华人民共和国教育部}
\ith{放弃入学}
\ist{入学前后}
\begin{qa}{E0343}{报到要带什么材料？准考证还需要吗？}
结论：推免生不走统考初试，没有准考证这回事；报到材料以录取通知书、身份证和双证原件为核心。国家规定，按国家招生规定录取的新生持录取通知书按学校有关要求和规定的期限到校办理入学手续，学校在报到时对入学资格进行初步审查（\nref{S0040} 第八条、第九条）。哈尔滨工业大学要求已录取推免生在入学报到时将身份证、学历学位证书原件提供给录取学院审核，报到前未取得本科毕业证书即取消录取资格（\nref{S0066}）。入学后3个月内还会做一轮复查：西北工业大学列明包括照片复查、证件复查、档案复查、资格复查、户籍复查（\nref{S0506}）；中央民族大学列出的复试资格审核材料含思想政治表现审查表原件、学历/学籍认证材料、考生诚信承诺书等（\nref{S0279}）。
\nref{S0040}
\nref{S0066}
\nref{S0506}
\nref{S0279}
\end{qa}

\iin{中华人民共和国教育部}
\ith{报到材料}
\ist{入学前后}
\begin{qa}{E0348}{学籍注册和录取备案分别是什么时间？}
结论：两者不是一回事，时点也差得很远。录取备案在推免系统里完成：推免生须在规定时间内通过系统答复是否接受“待录取”通知，待录取通知可拒绝多个但只能接受一个，接受后即成为该招生单位的待录取考生，不能再填报其他志愿（\nref{S0444}）；学校在系统备案后录取才生效——河北大学明确“未在教育部'推免服务系统'获得推免资格备案的考生，拟录取资格及其公示名单自动失效”（\nref{S0604}）。学籍则要等报到：学校在报到时对新生入学资格进行初步审查，审查合格办理入学手续、予以注册学籍（\nref{S0040} 第九条）；此后每学期开学都要按学校规定办理注册手续，未按规定缴纳学费或有其他不符合注册条件的，不予注册（第十二条）；学籍学历还须按相关规定完成电子注册（第三十五条）。
\nref{S0444}
\nref{S0604}
\nref{S0040}
\end{qa}

\iin{教育部学生服务与素质发展中心（学信网/研招网）}
\ith{学籍注册}
\ist{入学前后}
\section{延毕风险}
\begin{qa}{E0335}{哪些情况最容易延毕？怎么提前避坑？}
结论：延毕的成因集中在学业进度与培养环节两类，且大多可以提前识别。教育部要求培养单位抓住学位论文开题、中期考核、评阅、答辩、学位评定等关键环节，严格执行全方位全流程管理，并对不适合继续攻读学位的研究生“落实及早分流，加大分流力度”（\nref{S0432}）。制度上的硬约束是最长学习年限：《普通高等学校学生管理规定》（\nref{S0040}）要求除另有规定外，学生应在学校规定的最长学习年限（含休学和保留学籍）内完成学业（第二十五条），学业成绩未达到学校要求或未在规定学习年限内完成学业的，学校可予退学（第三十条）。中办国办《关于加快推动博士研究生教育高质量发展的意见》也提出强化分流退出和多向选择（\nref{S0459}）。
\nref{S0432}
\nref{S0040}
\nref{S0459}
\end{qa}

\iin{教育部学生服务与素质发展中心（学信网/研招网）}
\ith{延毕风险}
\ist{入学前后}
\begin{qa}{E0336}{硕士延毕常见原因是什么？博士呢？}
结论：诱因不同，学制与收费口径也不同。硕士延毕多与论文进度和培养环节不达标有关；博士延毕更多与课题不确定性、成果发表周期挂钩。可依据的硬数据是学制与最长年限：中国农业大学规定直博生学制5年、硕博连读生博士阶段学制3年、申请考核制博士生学制4年，最长在校学习年限为学制加2年（\nref{S0160}）；河北大学规定直博生学制5年，若因客观原因不能按期完成学业可申请延长学习年限，但在校最长学习年限不得超过8年（\nref{S0604}）；电子科技大学规定提前毕业或延期毕业的学生按实际学习年限收费（\nref{S0495}）；教育部办公厅也要求培养单位抓住开题、中期考核、评阅、答辩等关键环节，对不适合继续攻读学位的落实及早分流（\nref{S0432}）。

\par\vspace{2pt}\nobreak\begin{center}\footnotesize
\begin{tabularx}{\textwidth}{>{\raggedright\arraybackslash}X>{\raggedright\arraybackslash}X>{\raggedright\arraybackslash}X}
\toprule
\textbf{} & \textbf{硕士} & \textbf{博士} \\ \midrule
常见卡点 & 论文进度、答辩窗口 & 课题不确定性、成果要求 \\
关键环节 & 开题、答辩 & 中期考核、预答辩、评阅 \\
共同诱因 & 开题拖延、导师沟通断档、中期考核未过 & 同左 \\
\bottomrule
\end{tabularx}
\end{center}\vspace{2pt}\par
\nref{S0160}
\nref{S0604}
\nref{S0495}
\nref{S0432}
\end{qa}

\iin{中国农业大学}
\ith{延毕风险}
\ist{入学前后}
\begin{qa}{E0337}{直博读不下去怎么办？能转成硕士毕业吗？}
结论：制度上是支持的，但有学校层面的代价条款，且要走分流程序。中办国办《关于加快推动博士研究生教育高质量发展的意见》明确要求“强化分流退出和多向选择”（\nref{S0459}）；教育部办公厅要求培养单位对不适合继续攻读学位的研究生落实及早分流、加大分流力度（\nref{S0432}）。具体到操作与代价，郑州大学2027年办法写明“直博生如按硕士研究生毕业，须退还在学期间享受的学校普通助学金博士生与硕士生待遇的差额”（\nref{S0635}）。也就是说，转硕这条路是通的，但可能涉及待遇差额的退还，各校规定不同。

\begin{itemize}
  \item 时点：越早提越好，通常在博士资格考核或中期考核节点前
  \item 后果：可能需退还博士与硕士待遇差额
  \item 通道依据：及早分流、多向选择是主管部门明确要求的方向
\end{itemize}
\nref{S0459}
\nref{S0432}
\nref{S0635}
\end{qa}

\iin{教育部学生服务与素质发展中心（学信网/研招网）}
\ith{直博转硕}
\ist{入学前后}

\begin{qa}{E0350}{录取后要提前进实验室，是否等同于研究生学籍已经注册？}
不等同。提前进组是导师或课题组的安排，不会提前产生研究生学籍，也不能代替正式报到、资格复审和注册。开始前先向导师确认是否自愿、起止时间、住宿与安全责任、补贴、保险、实验室培训和数据保密要求；涉及实验操作应完成单位要求的安全培训并遵守本科生身份适用的权限。教育部规定新生入学后由学校按规定办理学籍注册，报到前仍须满足录取资格复核条件。若提前进组与毕业课程、论文或实习冲突，应先取得本科学校批准，避免影响毕业和推免资格。
\nref{S0045}
\nref{S0576}
\nref{S0762}
\qscope{操作原则适用于一般情形；具体流程、期限和便利措施以招生单位当年正式通知为准。}
\end{qa}
